\chapter{Similarity, Ratios, and Scale}
\label{ch:ch01-similarity-ratios}

\section{Why Size Matters}

Imagine two maps of the same city, one small enough to fit on a pocket card
and the other large enough to cover a classroom wall. The streets, rivers, and
landmarks occupy different physical distances on the two sheets of paper, yet
both maps represent the same place: the larger map has not created new
streets, and the smaller map has not removed any, since one is simply a
scaled copy of the other. This observation introduces one of the oldest and
most important ideas in mathematics, namely that some properties change when
size changes while others remain the same. A line segment that measures two
centimeters on one map may measure ten centimeters on another, and areas and
volumes change along with it, yet certain relationships survive: the ratio
between the lengths of two streets is unchanged, the angle at which two roads
meet is unchanged, and the overall shape of the city persists even as its
scale varies from map to map. Mathematics often begins by asking a
deceptively simple question --- what remains unchanged when everything else
changes? --- and the study of similarity is one of the oldest and most
productive answers to that question.

\section{Scale Transformations}

\begin{definition}
Let $P = (x,y)$ be a point of the Cartesian plane and let $k$ be a positive
real number. The \emph{scale transformation}, or \emph{dilation}, by factor
$k$ sends $P$ to the point $P' = (kx, ky)$.
\end{definition}

Every point of a figure moves proportionally under this transformation: if
$k > 1$ the figure expands away from the origin, and if $0 < k < 1$ the
figure contracts toward it. The map $(x,y) \mapsto (kx,ky)$ changes the size
of a figure while preserving its shape, and the question this chapter takes
up is which properties survive such a transformation and which do not.

\section{Similar Figures}

\begin{definition}
Two figures are \emph{similar} if one can be obtained from the other by some
combination of scaling, translation, reflection, and rotation. Informally,
similar figures share a common shape without necessarily sharing a common
size.
\end{definition}

Two squares of different sizes are similar, as are two circles of different
radii and two equilateral triangles of different side lengths, since in each
case one figure is obtained from the other by scaling alone. A square and a
rectangle whose sides are unequal, by contrast, are not similar, since no
combination of scaling, translation, reflection, and rotation carries one
onto the other. Similarity is therefore a statement about structure rather
than about measurement: it asks whether two figures agree in shape, leaving
their absolute size unspecified.

\section{Similar Triangles}

Triangles provide the most important examples of similarity encountered in
this book. Consider two triangles $\triangle ABC$ and $\triangle DEF$. If
their corresponding angles agree, so that $\angle A = \angle D$,
$\angle B = \angle E$, and $\angle C = \angle F$, the triangles are said to be
similar, written $\triangle ABC \sim \triangle DEF$. When two triangles are
similar, every pair of corresponding sides occurs in the same ratio,
\[
\frac{AB}{DE} = \frac{BC}{EF} = \frac{AC}{DF},
\]
and this common ratio is called the \emph{scale factor} relating the two
triangles.

\section{Derivation: Why Ratios Survive Scaling}

\begin{theorem}
Let $a$ and $b$ be the lengths of two segments, and suppose both are scaled
by the same factor $k > 0$, becoming $ka$ and $kb$. Then the ratio of the new
lengths equals the ratio of the original lengths: $\dfrac{ka}{kb} =
\dfrac{a}{b}$.
\end{theorem}

\begin{proof}
Since $k \neq 0$, it may be cancelled from numerator and denominator:
\[
\frac{ka}{kb} = \frac{k}{k}\cdot\frac{a}{b} = \frac{a}{b}.
\]
\end{proof}

The scale factor cancels because it appears identically in numerator and
denominator, and this is the fundamental reason ratios play such a central
role in geometry. A ratio ignores absolute size and records only structure:
whenever every length in a figure changes by the same factor, the ratios
between those lengths remain exactly as they were before the figure was
scaled.

\section{The Geometry of Invariance}

\begin{example}
Consider a right triangle with sides $3$, $4$, and $5$, and suppose every
side is doubled, producing a new triangle with sides $6$, $8$, and $10$. The
lengths, the area, and the perimeter of the new triangle all differ from
those of the original. Yet
\[
\frac{3}{5} = \frac{6}{10} \qquad\text{and}\qquad \frac{4}{5} = \frac{8}{10},
\]
so the ratios between corresponding sides remain exactly the same: only the
scale has changed, not the shape.
\end{example}

This observation lies at the foundation of everything that follows in this
book. The trigonometric functions to be introduced in later chapters arise
precisely because certain ratios of a right triangle depend only on the
triangle's angles and not on its size, and before sine, cosine, or tangent
can be defined in any meaningful way, it is necessary to first understand why
such size-independent quantities can exist at all. Similarity provides the
answer: because ratios are invariant under scaling, a ratio computed from any
triangle similar to a given one will agree with the ratio computed from the
original, which is exactly the property that will let the trigonometric
ratios of Chapter~3 depend on angle alone.

\section{Similarity Criteria}

Three standard criteria determine whether two triangles are similar without
requiring every angle and every side to be checked individually.

\begin{theorem}[Angle--Angle]
If two angles of one triangle equal two corresponding angles of another, the
two triangles are similar.
\end{theorem}
\begin{proof}
Since the angles of a triangle sum to a fixed total, agreement of two
corresponding angles forces agreement of the third as well, so all three
corresponding angles agree and the triangles are similar by definition.
\end{proof}

\begin{theorem}[Side--Angle--Side]
If two pairs of corresponding sides of two triangles are proportional and the
included angles are equal, the triangles are similar.
\end{theorem}

\begin{theorem}[Side--Side--Side]
If all three pairs of corresponding sides of two triangles occur in the same
ratio, $\dfrac{AB}{DE} = \dfrac{BC}{EF} = \dfrac{AC}{DF}$, the triangles are
similar.
\end{theorem}

\begin{algorithmproc}{How to Determine Whether Two Triangles Are Similar}
Given two triangles, first check whether two pairs of corresponding angles
are known and equal; if so, the Angle--Angle criterion applies and no further
check is needed. If instead two pairs of corresponding sides and their
included angle are known, compute the ratio of the two known sides in each
triangle; if the ratios agree and the included angles are equal, the
Side--Angle--Side criterion applies. If all three pairs of corresponding
sides are known, compute all three ratios; if the three ratios agree, the
Side--Side--Side criterion applies. If none of these conditions can be
verified from the given information, similarity cannot yet be concluded, and
more information about the triangles is required.
\end{algorithmproc}

\begin{example}
Determine whether a triangle with sides $3$, $4$, $5$ is similar to a
triangle with sides $6$, $8$, $10$. Computing the three ratios of
corresponding sides gives
\[
\frac{6}{3} = 2, \qquad \frac{8}{4} = 2, \qquad \frac{10}{5} = 2,
\]
and since all three ratios agree, the Side--Side--Side criterion applies:
the triangles are similar, with scale factor $k = 2$.
\end{example}

\section{Applications: Similarity as a Tool for Measurement}

Similarity allows unknown lengths to be determined without direct
measurement. Suppose a tree casts a shadow $12$ meters long, and at the same
moment a vertical meter stick casts a shadow $0.8$ meters long. Because
sunlight arrives at essentially parallel angles, the triangle formed by the
tree and its shadow is similar to the triangle formed by the meter stick and
its shadow, so the ratio of height to shadow length must agree between them:
\[
\frac{h}{12} = \frac{1}{0.8}.
\]
Solving for $h$ gives $h = 15$, so the tree is fifteen meters tall. No direct
measurement of the tree itself was required; similarity alone was enough to
convert an inaccessible length into an accessible one.

\begin{practicenow}
A flagpole casts a shadow $9$ meters long at the same moment that a
$1.5$-meter fence post casts a shadow $2$ meters long. Set up the
similar-triangle ratio implied by these measurements and find the height of
the flagpole.
\end{practicenow}

\section{Looking Ahead}

This chapter has introduced the idea that will support everything that
follows: a scale transformation changes lengths but preserves the ratios
between them, and those preserved quantities are the first examples of what
this book will repeatedly call invariants. The next chapter introduces
coordinates and vectors, and with them a second kind of transformation ---
rotation --- that changes orientation while preserving other quantities in
turn. Where similarity studies what survives a change of size, the chapters
ahead study what survives a change of orientation, and the path from
triangles to the trigonometric functions begins with the same simple
cancellation derived above: $ka/kb = a/b$. The scale factor vanishes, and
what remains is structure.

\section{Exercises}
\begin{exercise}
A triangle has sides $5$, $12$, $13$. A second triangle has sides $10$,
$24$, $26$. Determine whether the triangles are similar, and if so, state the
scale factor.
\end{exercise}

\begin{exercise}
Two triangles share an angle of $40^\circ$, and the sides adjacent to that
angle are in the ratio $3:5$ in both triangles. Which similarity criterion
applies, and why is that criterion sufficient without knowing the remaining
angles or sides?
\end{exercise}

\begin{exercise}
A person of height $1.7$ meters stands beside a building and casts a shadow
$2.1$ meters long. At the same moment, the building casts a shadow $34$
meters long. Find the height of the building, and explain in a sentence why
the assumption of parallel sunlight is essential to the argument.
\end{exercise}
